El proyecto \textbf{ELCO-Dealer} ha cumplido satisfactoriamente con el objetivo de diseñar e integrar un sistema automatizado complejo, aplicando de forma práctica los conocimientos de electrónica, mecánica y software embebido bajo la filosofía de una ``pequeña empresa de ingeniería''.

A través de este desarrollo, se ha demostrado que es posible automatizar un proceso cotidiano con un alto grado de fiabilidad mediante el uso de componentes accesibles como la \textbf{Arduino MEGA 2560 R3} y sensores de precisión \cite{arduino_mega}. La integración de soluciones innovadoras, como el \textbf{Slip Ring} para el giro infinito y el modelado paramétrico en \textbf{OnShape}, eleva el prototipo por encima de las soluciones comerciales básicas actuales.

Desde el punto de vista de la \textbf{viabilidad económica}, el análisis realizado muestra un camino claro hacia la rentabilidad: el paso de un coste de prototipado de aproximadamente 200€ a un coste optimizado de fabricación en masa de entre \textbf{35€ y 50€} por unidad \cite{bom}. Esto, unido a la estrategia de difusión en redes de código abierto (\textbf{GitHub e Instructables}), posiciona al ELCO-Dealer como un concepto sólido, escalable y con un gran potencial de mercado dentro del sector del ocio y los juegos de mesa .